\section{Results, plots and export}
\label{sec:results}

\subsection{The Results window}

Three tabs:
\begin{itemize}
  \item \textbf{Selected-time result}: the numerical table for the chosen
        time (magnitude row in photometry; per-wavelength rows in
        spectroscopy), showing \emph{every} quantity the engine produced:
        rates, S/N, $\sigma$(EW) [m\AA] for spectroscopy, the
        scintillation and ADC-quantization noise in electrons, pixels in
        the aperture, peak electrons/ADU, saturation cause, maximum
        unsaturated exposure, the standard and instrumental magnitudes and
        the sky brightness used. \textsf{Save result CSV} exports the full
        table.
  \item \textbf{Time series}: one row per time sample with UTC and local
        time, MJD, elevation, azimuth, \textbf{parallactic angle}, airmass
        (Pickering, refraction-corrected), S/N or required exposure,
        \textbf{sky surface brightness} actually used, band-effective
        extinction and its per-airmass coefficient, the
        absorption-adjusted apparent magnitude, saturation flag, peak
        electrons and maximum unsaturated exposure. This table is the CSV
        export.
  \item \textbf{Assumptions}: a plain-language statement of every physical
        assumption of the current run --- sky model and its components
        (including the SQM value in \code{sqm} mode), PSF model, the CMOS
        gain-table setting in use, the spectrograph geometry, noise terms,
        saturation policy, template transforms --- plus the selected-time
        altitude, airmass, parallactic angle, the \textbf{pixel scale in
        arcsec/pixel}, and sky brightness.
  \item \textbf{Outputs}: the selected-time results, in a tab of their own:
        S/N, scintillation noise in electrons and as a percentage of the
        source, ADC noise, peak-pixel and saturation numbers, predicted
        magnitudes, the unreddened $(B-V)_0$ and reddened $(B-V)$ colours,
        the $\sigma$(EW) line criterion, the stack plan and the
        differential-photometry precision. For spectroscopy it also lists
        the RV shift in pixels and km/s-per-pixel sampling, the
        \textbf{spectrum length in pixels} with a pass/fail against the
        detector $L$; the \textbf{visible diffraction orders} (slitless)
        with their start/end pixel from the zero order and a flag if an
        order runs off the chip; and the \textbf{Horne (1986) optimal
        extraction width} in arcsec and pixels.
\end{itemize}

For \textbf{photometry}, the \textbf{Selected-time result} tab keeps the
usual row for the entered magnitude, then a divider, then a greyed block of
the same calculation at the target magnitude $\pm1\ldots\pm7$, so the whole
brightness response is visible at a glance.

\subsection{Scientific plots}

The native Matplotlib window provides pan/zoom/save, manual axis limits
(time axes accept ISO timestamps), autoscale per panel, and in
spectroscopy a manual slider that walks the sampled times of the observing
interval, updating the spectrum panels and the markers. The sky-path panel
shows target, Sun (daytime only) and Moon (only above the horizon) in
altitude--azimuth (N=0$^\circ$, E=90$^\circ$).

\subsection{Practical notes}

\begin{itemize}
  \item The Moon's effect is included automatically in \code{ing} sky mode
        whenever it is above the horizon; within $\sim30^\circ$ of a bright
        Moon expect the sky several magnitudes brighter.
  \item In \code{fixed\_ab} mode the entered sky brightness is used as-is
        at every time --- use it when you have a measured value.
  \item Airmass and all results are deliberately unavailable below
        $5^\circ$ altitude.
  \item Save an \textbf{instrument profile} once your telescope/detector
        column is right; it restores automatically at start-up and travels
        with its QE/response files.
\end{itemize}
